\begin{abstract}
Reasoning language models produce long chains of thought, making
\emph{early exit}---stopping generation once the answer is settled---an
attractive way to cut inference cost. A popular family of methods
(e.g., Dynasor / Certaindex) periodically probes the partial trajectory for
its current answer and stops once probes \emph{agree}, betting that
intermediate consensus signals termination. We show this bet is unsafe, and
that the failure is not repaired by tuning the stopping rule within this family.
First, on \textsc{MATH}500 we document \emph{false consensus}: even a unanimous
window of recent probes is correct only $93.5\%$ of the time, and a naive
consensus stop loses $16.4$ percentage points (pp) of accuracy
($85.6\%\!\rightarrow\!69.2\%$) because continued reasoning frequently
overturns the intermediate answer. Second, we run a \emph{preregistered} search
over a seven-dimensional rule schema ($17{,}712$ candidate rules) on frozen
reasoning trajectories from two models and three math benchmarks, with
problem-grouped train/dev/test splits and stopping gates fixed in advance. On
the development set (held-out confirmation pending), \emph{no rule clears the
preregistered safety bar}: the safest stop still costs $1.85\pp$ per model.
Separately, under our dense-probe accounting, any rule with positive \emph{net}
token savings costs at least $4.87\pp$, and---counterintuitively---adaptive,
event-triggered probing is dominated by simple window consensus. We decompose
the cost into an intrinsic \emph{accuracy tax} (independent of probing) and a
\emph{probe tax} (a function of probe density), and argue only the accuracy tax
is probe-invariant. Our results suggest that safe early exit for reasoning
models likely requires a different termination signal, not a better-tuned
confidence threshold on probe agreement.
\end{abstract}
